\chapter{Реализация}\label{ch:impl}

Настоящая глава описывает реализацию предложенной в главе~3 схемы
хранения в кодовой базе Otterbrix. Основная цель раздела~--- показать,
как архитектурные решения отображаются на конкретные компоненты
кода, и какие структурные изменения потребовалось внести в
различные подсистемы.

\section{Обзор затронутых компонентов}

Реализация задействует следующие ключевые подсистемы Otterbrix:

\begin{itemize}
    \item \texttt{components/catalog/}~--- метаданные таблиц.
    Введён новый класс \texttt{computed\_schema}, отвечающий за
    учёт зарегистрированных полей и их типов.
    \item \texttt{components/physical\_plan/operators/}~---
    физические операторы. Введены четыре новых оператора:
    \texttt{scan\_computing\_table} (гибридное сканирование),
    \texttt{operator\_insert\_computing},
    \texttt{operator\_delete\_computing},
    \texttt{operator\_update\_computing} (multi-collection DML под
    одной транзакцией).
    \item \texttt{services/collection/}~--- исполнитель планов.
    Цикл DML расширен поддержкой ``вычисляемых''
    DML-операторов: \texttt{intercept\_dml\_io\_} получил отдельные
    ветви для \texttt{insert\_computing},
    \texttt{remove\_computing} и \texttt{update\_computing},
    которые ведут многотабличный жизненный цикл транзакции (запись
    в N+1 таблиц, WAL physical, общий commit с одним
    $\mathit{commit\_id}$, общий COMMIT marker).
    \item \texttt{services/dispatcher/}~--- центральная точка
    маршрутизации запросов. Pre-flight для \textsc{insert} в
    вычисляемую таблицу создаёт недостающие side-таблицы и
    проставляет на узел плана список нужных side для последующего
    исполнителя.
    \item \texttt{services/disk/}~--- менеджер дискового
    хранилища. Добавлена поддержка создания, удаления и сканирования
    side-таблиц (как обычных регулярных таблиц).
    \item \texttt{components/sql/transformer/}~--- SQL-трансформер.
    Добавлена поддержка синтаксиса \texttt{CREATE TABLE db.t~();}
    (вычисляемая схема), \texttt{WITH (pinned\_columns~=~\ldots)}
    (список pinned-полей) и \texttt{INSERT \ldots\ VALUES (json('\ldots'))}.
    \item \texttt{components/logical\_plan/}~--- узлы логического
    плана. На \texttt{node\_create\_collection\_t} добавлен флаг
    динамической схемы и опциональный список pinned-полей; на
    \texttt{node\_insert\_t}, \texttt{node\_delete\_t} и
    \texttt{node\_update\_t}~--- метаданные \texttt{computing\_sides}
    для маршрутизации в специализированные операторы.
    \item \texttt{components/table/}~--- низкоуровневое колоночное
    хранилище. Без существенных модификаций~--- side-таблицы
    хранятся как обычные регулярные таблицы.
    \item \texttt{components/vector/}~--- операции над векторами
    значений. Добавлены placeholder-векторы (без выделения буфера
    для непроецированных колонок) для поддержки column~pruning в
    \texttt{scan\_computing\_table}.
\end{itemize}

\section{Класс представления вычисляемой схемы}

Класс \texttt{computed\_schema}
(\texttt{components/catalog/computed\_schema.hpp}/.cpp)
инкапсулирует учёт метаданных таблицы с вычисляемой схемой.

\subsection{Поля состояния}

Класс \texttt{computed\_schema} хранит следующее состояние:

\begin{itemize}
    \item отображение имени поля в перечень типов, в которых оно
    встречалось;
    \item упорядоченный список пар <<имя поля~--- тип>> в порядке
    их первого появления;
    \item множество имён полей, объявленных pinned при
    \texttt{CREATE TABLE} (если таковые есть);
    \item флаг \texttt{is\_sparse}, указывающий, является ли
    таблица таблицей с вычисляемой (разреженной) схемой; для обычных
    регулярных таблиц этот флаг равен \texttt{false} и логика
    side-таблиц отключена.
\end{itemize}

Динамическое отслеживание статистики заполненности и счётчики
non-null убраны вместе с механизмом автоматического продвижения
(см.~раздел ``Переход между режимами'' в главе~3).

Центральное понятие~--- имя side-таблицы, формируемое статической
функцией \texttt{side\_table\_name(main, field, type)} по правилу
\texttt{\_dyn\_<main>\_\_<field>\_\_<type\_enum>}. Префикс
\texttt{\_dyn\_} зарезервирован для служебных таблиц и гарантирует,
что автоматически порождённые имена не конфликтуют с
пользовательскими. Тип в суффиксе необходим потому, что одно и то
же логическое поле может существовать в нескольких типах: если
первый \textsc{insert} вставил \texttt{\{"price":~100\}}, а
следующий~--- \texttt{\{"price":~"free"\}}, в схеме появляются две
отдельные side-таблицы \texttt{\_dyn\_t\_\_price\_\_INT64} и
\texttt{\_dyn\_t\_\_price\_\_STRING}.

\subsection{Центральные операции}

\begin{description}
    \item[\texttt{append(field\_name, type)}] Регистрирует новую
    пару (field, type) в схеме, если её ещё нет. Идемпотентна.
    \item[\texttt{latest\_types\_struct()}] Возвращает текущую
    виртуальную схему таблицы в виде структурного типа: каждая
    зарегистрированная пара (field, type) фигурирует как поле.
    Используется валидатором для проверки запросов на
    soundness относительно текущей схемы.
    \item[\texttt{side\_table\_name(main, field, type)}]
    Генерирует имя side-таблицы. Статическая функция; не
    обращается к состоянию объекта.
\end{description}

\subsection{Потокобезопасность}

\texttt{computed\_schema} синхронизируется внешним блокирующим
мьютексом на уровне каталога. Все изменения схемы
(\texttt{append}) выполняются под мьютексом каталога во время
pre-flight \textsc{insert}. Поскольку изменения~--- лишь добавление
новых пар (field, type), а не модификация существующих, конкуренция
на мьютексе минимальна. Чтения схемы (\texttt{latest\_types\_struct})
требуются только при валидации запросов и физически не конфликтуют
ни с какими записями.

\section{Диспетчер: pre-flight для DML}

Обработчик каждого DML-узла (\textsc{insert}, \textsc{delete},
\textsc{update}) в диспетчере, обнаружив таблицу с вычисляемой
схемой (\texttt{catalog\_.table\_computes(id)}), выполняет короткую
pre-flight-фазу под мьютексом диспетчера и далее передаёт работу в
исполнитель:

\begin{itemize}
    \item Для \textsc{insert}: проход по колонкам входного чанка с
    регистрацией пары (поле, тип) в \texttt{computed\_schema} и
    созданием соответствующих side-таблиц для тех пар, что
    впервые увидены (атомарно через
    \texttt{manager\_disk\_t::create\_storage\_with\_columns},
    \texttt{register\_collection},
    \texttt{catalog\_.create\_table}). Список созданных и уже
    существующих side-таблиц проставляется на узел
    \texttt{node\_insert\_t} в виде поля
    \texttt{computing\_sides}.
    \item Для \textsc{delete} и \textsc{update}: текущая схема
    таблицы (полный список side) проставляется на соответствующий
    узел~--- без DDL-операций, поскольку DML только модифицирует
    существующие записи.
\end{itemize}

После pre-flight план отправляется в исполнитель стандартным
вызовом \texttt{execute\_plan\_impl}~--- ровно тот же путь, что и
для обычной DML на регулярной таблице. Мьютекс диспетчера
освобождается сразу после успешной маршрутизации, и продолжительная
DML-работа над main и side-таблицами выполняется на исполнителе
без удержания глобального мьютекса.

\section{Исполнитель: операторы вычисляемой DML}

Главное архитектурное изменение реализации~--- перенос всей
многотабличной DML-работы из диспетчера в специализированные
физические операторы в исполнителе. Это позволяет проводить запись
в main и все нужные side-таблицы под одной транзакцией с одним
$\mathit{txn\_id}$ и одним $\mathit{commit\_id}$, что даёт
real-MVCC-атомарность.

\subsection{operator\_insert\_computing}

Минималистичный read\_write-оператор: хранит имя main-таблицы и
список side (\texttt{computing\_sides}), полученный с
\texttt{node\_insert\_t}. В \texttt{on\_execute\_impl} переносит
вывод левого ребёнка (исходный \texttt{data\_chunk\_t} от
\texttt{operator\_data}) в собственный \texttt{output\_} и
переходит в режим \texttt{async\_wait}. Вся работа происходит в
ветке \texttt{intercept\_dml\_io\_::case insert\_computing} в
исполнителе:

\begin{enumerate}
    \item \texttt{storage\_total\_rows} на main~--- получение
    диапазона новых \texttt{\_id}.
    \item Разбиение пользовательского чанка: pinned-колонки и
    \texttt{\_id} формируют чанк для main, каждое sparse-поле
    формирует двухколоночный (\texttt{\_id}, value) чанк для своей
    side-таблицы (только ненулевые строки).
    \item \texttt{storage\_append} на каждую side-таблицу с
    текущим $\mathit{txn\_id}$; затем \texttt{storage\_append} на
    main с тем же $\mathit{txn\_id}$.
    \item WAL physical записи для каждой затронутой таблицы.
    \item \texttt{txn\_manager.commit} $\to$
    $\mathit{commit\_id}$;
    \texttt{storage\_commit\_append} для каждой таблицы с этим
    $\mathit{commit\_id}$.
    \item Один WAL COMMIT marker для всей транзакции.
\end{enumerate}

При ошибке на любом шаге выполняется
\texttt{storage\_revert\_append} для уже записанных таблиц и
\texttt{txn\_manager.abort}. Поскольку $\mathit{commit\_id}$ не
был выставлен, эти записи никому невидимы. Side-таблицы могут
содержать ``висящие'' уже-помеченные-удалёнными записи; они
безопасно отчищаются последующим garbage collection.

\subsection{operator\_delete\_computing}

В отличие от обычного \texttt{operator\_delete}, который работает
с одной таблицей, computing-вариант знает обо всех side-таблицах
схемы. Левое поддерево оператора~---
\texttt{scan\_computing\_table~+~operator\_match}: оно отдаёт
чанки с логическими \texttt{\_id} в \texttt{chunk.row\_ids} для
тех строк, что удовлетворяют WHERE.

В ветке \texttt{intercept\_dml\_io\_::case remove\_computing}:

\begin{enumerate}
    \item Собирается множество удаляемых \texttt{\_id}.
    \item Для каждой целевой таблицы (main и каждая side) делается
    \texttt{storage\_scan\_batched} с подбором физических позиций
    тех строк, чьё значение колонки \texttt{\_id} принадлежит
    собранному множеству. Это эквивалентно поиску по индексу
    \texttt{\_id}, если такой индекс есть, или линейному скану с
    pruning по zone~maps в противном случае.
    \item \texttt{storage\_delete\_rows} с найденными физическими
    позициями и текущим $\mathit{txn\_id}$ на каждую таблицу. Main
    помечается \emph{первой}, side~--- после, согласно инварианту
    ``main определяет видимость, удаление сначала отменяет видимость
    в main, потом чистит side''.
    \item WAL physical\_delete, \texttt{txn\_manager.commit},
    \texttt{storage\_commit\_delete} на каждую таблицу с общим
    $\mathit{commit\_id}$, и единственный WAL COMMIT marker.
\end{enumerate}

\subsection{operator\_update\_computing}

Реализует ``re-alloc с полным копированием'': для каждой
матчащейся строки выделяется новый \texttt{\_id}, под этот новый
\texttt{\_id} записываются все поля строки (с применением
SET-выражений), и старая строка удаляется. Структурно объединяет
ветви \texttt{insert\_computing} и \texttt{remove\_computing}.

Левое поддерево, аналогично DELETE, состоит из
\texttt{scan\_computing\_table~+~operator\_match}, но проектирует
\emph{все} колонки виртуальной схемы (не только row\_ids), чтобы
получить полные значения матчащихся строк. SET-выражения
применяются in-place к выходному чанку через
\texttt{update\_expr\_t::execute}, после чего:

\begin{enumerate}
    \item \texttt{storage\_total\_rows} $\to$ новый диапазон
    \texttt{\_id}.
    \item Сборка новых per-side и main-чанков, аналогично
    \texttt{insert\_computing}.
    \item \texttt{storage\_append} в side, затем в main.
    \item Сбор физических позиций старых \texttt{\_id} в каждой
    таблице (сканированием с фильтрацией).
    \item \texttt{storage\_delete\_rows} на тех же таблицах под
    тем же $\mathit{txn\_id}$.
    \item WAL physical для всех вставок и удалений,
    \texttt{txn\_manager.commit},
    \texttt{storage\_commit\_append}+\texttt{storage\_commit\_delete}
    на каждую таблицу с одним $\mathit{commit\_id}$, общий WAL
    COMMIT marker.
\end{enumerate}

Под одним $\mathit{commit\_id}$ становятся видны одновременно и
``новая'' строка, и снятие видимости со ``старой''~--- результат
наблюдаем как атомарная подмена.

\section{Оператор scan\_computing\_table}

Гибридное сканирование для \textsc{select} реализовано отдельным
физическим оператором
\texttt{components/physical\_plan/operators/scan/scan\_computing\_table.\{hpp,cpp\}}.

Оператор принимает на вход:
\begin{itemize}
    \item имя main-таблицы;
    \item список side (имя~--- тип значения), параллельный позиции
    в виртуальной схеме (без main);
    \item номер pinned-колонок main (доступен через
    \texttt{computed\_schema});
    \item список запрошенных позиций виртуальной схемы
    (\texttt{projected\_cols}~--- стандартный механизм
    column~pruning).
\end{itemize}

Алгоритм \texttt{await\_async\_and\_resume}:

\begin{enumerate}
    \item \emph{Параллельный запуск сканов.} Отправляются
    одновременно асинхронные сообщения
    \texttt{storage\_scan\_batched} к менеджеру диска для main и
    для каждой запрошенной side-таблицы. \texttt{co\_await}
    собирает все ответы; благодаря акторной модели Otterbrix
    физическое чтение разных таблиц может идти на разных
    исполнителях диска параллельно.
    \item \emph{Сканирование main.} Выходные чанки main отдают
    значения pinned-колонок напрямую (они физически здесь же), а
    также набор живых \texttt{\_id} (записываются в
    \texttt{chunk.row\_ids} выходных чанков как канал
    мета-идентификаторов, доступный последующим операторам:
    operator\_match, operator\_delete\_computing,
    operator\_update\_computing).
    \item \emph{Построение хеш-таблицы
    \texttt{\_id}~$\to$~позиция-в-выходе.} Один проход по
    собранным \texttt{\_id} строит словарь.
    \item \emph{Сканирование side с подстановкой.} Для каждой
    запрошенной sparse-side: проход по её чанкам, для каждой
    строки $(\_id, value)$ ищется позиция-в-выходе в хеш-таблице,
    \texttt{value} записывается в соответствующую колонку
    выходного чанка по этой позиции. Строки без записи в side
    остаются NULL (выходные векторы инициализированы как
    ``все~NULL'').
\end{enumerate}

Сложность по времени: $O(N + \sum_i n_i)$, где $N$~--- число
живых строк main, $n_i$~--- число строк в $i$-й запрошенной
side-таблице. Сложность по памяти: $O(N)$ для хеш-таблицы плюс
$O(N \cdot k)$ для выходных буферов, где $k$~--- число запрошенных
колонок (стандартно). Side-таблицы, не упомянутые в запросе, не
сканируются вовсе.

\section{SQL-трансформер}

\subsection{Создание таблицы с вычисляемой схемой}

Для создания таблицы с вычисляемой схемой используется синтаксис с
пустым списком колонок. SQL-трансформер
(\texttt{components/sql/transformer/impl/transform\_table.cpp})
распознаёт \texttt{CreateStmt} с пустым \texttt{tableElts} и создаёт
\texttt{node\_create\_collection\_t} с флагом
\texttt{is\_dynamic\_schema\_table=true}. Опционально через
синтаксис WITH передаётся список pinned-полей~--- они становятся
физическими колонками main-таблицы наряду со служебным
\texttt{\_id}:
\begin{verbatim}
CREATE TABLE db.t ();
CREATE TABLE db.t ()
  WITH (pinned_columns = 'commit.rev,commit.operation');
\end{verbatim}
Все остальные поля, появившиеся в \textsc{insert}, автоматически
становятся sparse и получают собственные side-таблицы. Дальнейшее
изменение списка pinned-полей не предусмотрено~--- статус поля
фиксируется при \texttt{CREATE TABLE} (см.~раздел
``Переход между режимами'' в главе~3).

\subsection{Вставка JSON-документов одним выражением}

Этот синтаксис предоставляет удобный способ вставки JSON-документов
напрямую из SQL. SQL-трансформер детектирует конструкцию
\texttt{json('\ldots')} как вызов функции (FuncCall) с одним
строковым аргументом в контексте VALUES-строки без указания списка
колонок \textsc{insert}. При обнаружении такой конструкции
активируется путь \texttt{build\_json\_insert}:
\begin{enumerate}
    \item Каждая VALUES-строка разбирается как JSON-документ через
    библиотеку Boost.JSON.
    \item JSON рекурсивно <<расплющивается>>: вложенные объекты
    порождают пути с точечной нотацией (\texttt{user.profile.age}),
    массивы сериализуются обратно в JSON-строки, null-значения
    пропускаются.
    \item По всему батчу строится объединение путей. Для каждого
    пути определяется тип на основании первого ненулевого значения.
    \item Создаётся \texttt{data\_chunk\_t} с колонками по числу
    уникальных путей; для каждой строки проставляются значения
    найденных в ней полей.
    \item Формируется \texttt{node\_insert\_t}, который далее идёт
    в диспетчер и попадает в стандартный путь \textsc{insert} для
    вычисляемой схемы.
\end{enumerate}

\section{Модификации подсистемы хранения}

Side-таблицы хранятся как обычные регулярные двухколоночные
таблицы; все стандартные операции \texttt{storage\_append},
\texttt{storage\_delete\_rows}, \texttt{storage\_scan\_batched},
\texttt{storage\_commit\_*} работают с ними без модификаций.
Никаких специальных низкоуровневых структур для разреженного
хранения не требуется.

Единственная заметная правка~--- это поддержка ``placeholder''
векторов в \texttt{components/vector/data\_chunk.cpp}: для
column~pruning в \texttt{scan\_computing\_table} непроецированные
колонки виртуальной схемы получают \texttt{vector\_t} без
выделенного буфера, что экономит память и аллокатор. Операторы
вверх по плану (match, group, sort, select) корректно пропускают
такие placeholder-векторы благодаря проверке на null-buffer.

\section{Параллельность и корректность}

Корректность параллельной работы с таблицами вычисляемой схемы
подробно разобрана в разделе~\ref{sec:parallel} главы~3 на уровне
архитектуры. Здесь резюмируются её реализационные точки опоры в
кодовой базе.

\begin{itemize}
    \item \textbf{Атомарность ``side $\to$ main''.} Порядок записи
    реализован в \texttt{intercept\_dml\_io\_::case
    insert\_computing}: сначала
    \texttt{storage\_append} на все side-таблицы под общим
    $\mathit{txn\_id}$, затем \texttt{storage\_append} на main с
    тем же $\mathit{txn\_id}$. \texttt{storage\_commit\_append}
    с общим $\mathit{commit\_id}$ применяется ко всем таблицам в
    одной фазе после \texttt{txn\_manager.commit}.
    \item \textbf{Симметричный порядок для DELETE/UPDATE.} В
    \texttt{intercept\_dml\_io\_::case remove\_computing} удаление
    выполняется сначала на main (\texttt{storage\_delete\_rows}),
    затем на side. То же в \texttt{update\_computing}: вставка
    новой версии под новый \texttt{\_id} (side $\to$ main),
    удаление старой версии (main $\to$ side), всё под общим
    $\mathit{txn\_id}$ и $\mathit{commit\_id}$.
    \item \textbf{Атомарность выделения \texttt{\_id}.} Запрос
    \texttt{storage\_total\_rows} обрабатывается актёрным
    менеджером диска последовательно, что гарантирует
    непересекающиеся диапазоны \texttt{start\_id} для параллельных
    \textsc{insert}.
    \item \textbf{MVCC-снимок на каждый \textsc{select}.}
    \texttt{scan\_computing\_table} использует
    \texttt{ctx->txn.start\_time} при обращении к
    \texttt{storage\_scan\_batched}; menedger диска фильтрует
    видимые записи по \texttt{start\_time} читателя как для main,
    так и для side. Согласованность чтения нескольких таблиц одной
    транзакции обеспечивается тем, что
    $\mathit{start\_time}$ один и тот же для всех сканов.
    \item \textbf{Независимость от хеш-маршрутизации.} Текущая
    реализация Otterbrix маршрутизирует запросы к одной коллекции
    на один акторный исполнитель, что само по себе сериализует
    операции с этой коллекцией. Все приведённые выше аргументы о
    параллельности \emph{не зависят} от этой сериализации; они
    опираются только на MVCC-механизмы и инвариант ``main $\to$
    side''. Это специально проверено: ослабление маршрутизации не
    требует модификаций в коде вычисляемой схемы.
\end{itemize}

Совокупно система обеспечивает согласованность в духе
snapshot~isolation: каждая транзакция видит консистентное состояние
таблицы~--- включая sparse-поля,~--- независимо от параллельных
операций.

\section{Тестовое покрытие}

Для реализации разработано обширное тестовое покрытие, включающее:

\begin{itemize}
    \item \textbf{Модульные тесты}~--- проверка корректности
    \texttt{computed\_schema} в изоляции; распознавание SQL-
    конструкций \texttt{CREATE TABLE db.t~()} и
    \texttt{INSERT \ldots\ VALUES (json(\ldots))}.
    \item \textbf{Интеграционные тесты}~--- более 30~сценариев:
    базовый \textsc{insert}/\textsc{select} с автоматической
    эволюцией схемы; эволюция между батчами; разреженное хранение;
    продвижение с миграцией данных; pinned-колонки; запросы,
    смешивающие плотные и разреженные поля; граничные случаи.
    \item \textbf{Тесты производительности}~--- сравнение трёх
    конфигураций Otterbrix между собой и с DuckDB на стандартном
    датасете JSONBench.
\end{itemize}

Полный прогон всех тестов в конфигурации Debug занимает около
50~секунд и насчитывает свыше 470~тестовых сценариев. Все тесты на
момент завершения работы проходят успешно.

\chapterconclusion

Реализация предложенной архитектуры потребовала введения нового
класса метаданных \texttt{computed\_schema}, четырёх новых
физических операторов (\texttt{scan\_computing\_table} для
гибридного сканирования, \texttt{operator\_insert\_computing} /
\texttt{operator\_delete\_computing} /
\texttt{operator\_update\_computing} для multi-collection DML под
одной транзакцией), pre-flight-фазы в диспетчере для маршрутизации
DML-узлов в специализированные операторы и DDL side-таблиц,
расширения SQL-трансформера поддержкой динамического
\texttt{CREATE TABLE} и \texttt{INSERT \ldots\ VALUES
(json(\ldots))}, и минимальных правок низкоуровневых подсистем
хранения. Интеграция с существующей акторной моделью параллелизма
не потребовала введения новых механизмов синхронизации; согласованность
с MVCC и WAL обеспечивается выбором правильного порядка
``side $\to$ main'' для записей и ``main $\to$ side'' для удалений
под общим $\mathit{txn\_id}$ и единым $\mathit{commit\_id}$.
